\section{Two-body decay and laboratory acceptance}
\label{v571:kinematics}

The beam energy controls the angular separation of the electron and positron
from a boosted dark photon.  We first consider a prompt $A'\to e^+e^-$ decay
at the target, with the parent directed along the beam axis $z$.  In units with
$c=1$, let $m$ denote the parent mass and write its laboratory energy as
$E_A=xE_{\rm beam}$, where $x$ is the fraction of the beam energy carried by
the parent.  At fixed $m$ and $x$, the boost and daughter rest-frame momentum are
\begin{equation}
  \gamma=\frac{E_A}{m},\qquad
  \beta=\sqrt{1-\frac{m^2}{E_A^2}},\qquad
  E_* = \frac{m}{2},\qquad
  p_* = \sqrt{\frac{m^2}{4}-m_e^2},
  \label{v571:boost}
\end{equation}
with $2m_e\leq m<E_A$.  Define $c_*=\cos\theta_*$ and let $\phi_*$ be the
rest-frame azimuth measured from the horizontal $x$ axis.  The two rest-frame
four-vectors and their laboratory components are
\begin{equation}
  p_{\pm}^{*\mu}
    =\left(\frac m2,\ \pm p_*\sqrt{1-c_*^2}\cos\phi_*,\
       \pm p_*\sqrt{1-c_*^2}\sin\phi_*,\ \pm p_*c_*\right),
       \label{v571:rest-vectors}
\end{equation}
\begin{align}
  E_\pm&=\gamma\left(\frac m2\pm\beta p_*c_*\right),&
  p_{z,\pm}&=\gamma\left(\frac{\beta m}{2}\pm p_*c_*\right),
       \label{v571:lab-vectors}\\
  p_{x,\pm}&=\pm p_*\sqrt{1-c_*^2}\cos\phi_*,&
  p_{y,\pm}&=\pm p_*\sqrt{1-c_*^2}\sin\phi_* .
       \label{v571:transverse-vectors}
\end{align}
The signs label the opposite daughter directions.  For a forward daughter,
its polar angle and vertical projected angle are respectively
\begin{equation}
  \tan\theta_\pm=\frac{p_*\sqrt{1-c_*^2}}{p_{z,\pm}},\qquad
  \tan\theta_{y,\pm}=\frac{p_{y,\pm}}{p_{z,\pm}}.
  \label{v571:angles}
\end{equation}
The energy fractions are $E_\pm/E_A=(1\pm\beta\rho c_*)/2$, where
$\rho=\sqrt{1-4m_e^2/m^2}$.  They approach the momentum fractions only when
the daughters are ultrarelativistic.  Thus fixing $x$ fixes the pair energy,
while the decay angle still determines how that energy is shared.

\subsection{Equal-sharing mass scales}
\label{v571:equal-sharing}
For $c_*=0$, both daughters have energy $E_A/2$ and the same polar angle
$\theta$; their opening angle is $2\theta$.  The exact invariant-mass relation is
\begin{equation}
  m^2=E_A^2\sin^2\theta+4m_e^2\cos^2\theta
  \quad\longrightarrow\quad
  m\simeq xE_{\rm beam}\sin\theta
  \simeq xE_{\rm beam}\theta .
  \label{v571:equal-sharing-mass}
\end{equation}
The first approximation neglects the electron mass and the second assumes a
small angle.  In the vertical decay plane, $|\sin\phi_*|=1$, this gives the
mass scale associated with a specified vertical angle.  For a general decay,
the small-angle, relativistic relation is instead
$m^2\simeq E_+E_-\theta_{+-}^2$, with
$\theta_{+-}$ the daughter opening angle.  Unequal sharing therefore changes
the correspondence between a mass and either single-track angle.

A vertical beam gap excludes tracks close to the beam \emph{plane}; a polar
cone excludes tracks close to the beam \emph{axis}.  Their acceptances differ
because $|\tan\theta_y|=\tan\theta\,|\sin\phi_*|$.  In particular, an outer
vertical edge gives an upper equal-sharing mass scale only in the specified
vertical plane.  With unrestricted azimuth it does not, by itself, give an
upper mass cutoff.  Horizontal sensor boundaries and magnetic bending must
also be considered before interpreting an endpoint as detector acceptance.

\subsection{Normalized decay fraction through the tracker}
\label{v571:acceptance}
For the conditional angular calculation, we use a transverse-vector decay
distribution in the relativistic-daughter approximation, with uniform
azimuth:
\begin{equation}
  w(c_*)=\frac38(1+c_*^2),\qquad
  \mathcal A(m,x)=\int_{-1}^{1}\!dc_*\,w(c_*)
    \int_0^{2\pi}\!\frac{d\phi_*}{2\pi}\,
    \mathcal I_{\rm geom}(p_+,p_-).
  \label{v571:acceptance-integral}
\end{equation}
Here $\mathcal I_{\rm geom}$ is one when both forward daughters satisfy the
chosen geometrical requirements and zero otherwise.  The integral is
normalized to all decays with that specified polarization, parent energy and
direction.  Finite $m_e$ is retained in the kinematics above; the adopted
angular weight is not an exact threshold or production-polarization model.

For an active sensor rectangle centered at $\boldsymbol r_j$, let
$\boldsymbol n_j$ be its normal, $\boldsymbol a_j,\boldsymbol b_j$ its
orthogonal in-plane unit vectors, and $h_{a,j},h_{b,j}$ its active half-widths.
For the zero-field reference, a straight daughter ray
$\boldsymbol r(t)=\boldsymbol r_0+t\boldsymbol u$,
with $\boldsymbol u=\boldsymbol p/|\boldsymbol p|$, crosses the sensor plane at
\begin{align}
  t_j&=\frac{(\boldsymbol r_j-\boldsymbol r_0)\cdot\boldsymbol n_j}
                {\boldsymbol u\cdot\boldsymbol n_j},&
  \Delta\boldsymbol r_j&=\boldsymbol r_0+t_j\boldsymbol u-\boldsymbol r_j,
  \nonumber\\
  |\Delta\boldsymbol r_j\cdot\boldsymbol a_j|&\leq h_{a,j},&
  |\Delta\boldsymbol r_j\cdot\boldsymbol b_j|&\leq h_{b,j}.
  \label{v571:sensor-intersection}
\end{align}
A geometrical sensor hit requires $t_j>0$, a nonparallel ray, and both bounds.
For station $k$ in a specified detector half, define binary flags
$h_{k,a}$ and $h_{k,s}$ for its axial and stereo views.  Each flag is one if
the ray crosses any active rectangle belonging to that view.  Where hole and
slot sensors are alternatives in one view, their rectangles form a union:
crossing either is sufficient, and that view is counted at most once.
The geometrical counts for one daughter are then
\begin{equation}
  N_{\rm 2D}=\sum_k(h_{k,a}+h_{k,s}),\qquad
  N_{\rm 3D}=\sum_k h_{k,a}h_{k,s}.
  \label{v571:hit-counts}
\end{equation}
Thus a 2D plane hit requires one view, whereas a paired 3D space point
requires both views at the same station and in the same detector half.
These counts are not interchangeable when only one view is crossed:
$N_{\rm 2D}=2N_{\rm 3D}$ holds only if every crossed view has its partner.
The run- and charge-dependent thresholds, including any required inner-plane
hits, are specified in the selection table below.  They enter
$\mathcal I_{\rm geom}$ separately for the electron and positron.
Equation~\eqref{v571:sensor-intersection} gives the analytic crossing for the
zero-field reference. With a magnetic field, the curved daughter trajectory
determines the sensor crossings before the same view and station flags are
evaluated. These counts omit scattering and hit inefficiency.

\begin{samepage}
\subsubsection{Magnetic propagation}
\label{v571:magnetic-propagation}
Let $s$ denote path length, $q=+1$ for the positron and $q=-1$ for the
electron, and $\boldsymbol B(\boldsymbol r)$ the static field in the detector
coordinate frame. For momentum in MeV, position in mm and field in tesla,
the transport equations are
\begin{equation}
 \begin{gathered}
  \frac{d\boldsymbol r}{ds}=\boldsymbol u
      =\frac{\boldsymbol p}{|\boldsymbol p|},\qquad
  \frac{d\boldsymbol p}{ds}
      =\kappa q\,\boldsymbol u\times\boldsymbol B(\boldsymbol r),\\[3pt]
  \kappa=0.299792458\ \frac{\mathrm{MeV}}{\mathrm{T\,mm}}.
 \end{gathered}
  \label{v571:magnetic-equations}
\end{equation}
\end{samepage}
The magnetic force is perpendicular to the momentum, so
$d|\boldsymbol p|^2/ds=0$: each daughter's momentum magnitude and energy
remain constant in this transport model. The field changes the trajectory
and the directions at the sensors. The mass labeling the acceptance curves
is the invariant mass of the two momenta at the production vertex, rather
than a mass recomputed from their different local directions downstream.
Energy loss and material scattering require additional physics.

The numerical trajectory crosses sensor plane $j$ at a forward path length
$s_j>0$ satisfying
\begin{equation}
  \bigl[\boldsymbol r(s_j)-\boldsymbol r_j\bigr]\cdot
  \boldsymbol n_j=0.
  \label{v571:curved-crossing}
\end{equation}
The in-plane half-width bounds of Eq.~\eqref{v571:sensor-intersection}
are then tested at $\boldsymbol r(s_j)$. This prescription follows each
charge through the same field and retains the complete finite sensor
geometry. Setting $\boldsymbol B=0$ recovers the straight-ray calculation.
The field map, integration method and numerical checks are specified with
the detector configuration below.

\subsubsection{Projected angular reference}

For an illustrative projection, consider symmetric active
vertical intervals $y_{\min,j}\leq |y|\leq y_{\max,j}$ at longitudinal
positions $z_j$, measured from a target at $z_0$.  A straight trajectory from
the target has $y_j=(z_j-z_0)p_y/p_z$.  Requiring each daughter to cross every
chosen interval reduces to a constant slope band
\begin{equation}
  L\equiv\max_j\frac{y_{\min,j}}{z_j-z_0}
    \leq\left|\frac{p_y}{p_z}\right|\leq
  U\equiv\min_j\frac{y_{\max,j}}{z_j-z_0},
  \label{v571:projected-band}
\end{equation}
provided $0\leq L<U$.  This band is a reference for understanding the angular
scales, rather than a replacement for the finite sensor calculation.

The azimuthal integral can be performed explicitly.  At fixed $c_*$, set
$p_T=p_*\sqrt{1-c_*^2}$ and, for $p_T>0$ and $p_{z,+},p_{z,-}>0$, define
\begin{equation}
  q_{\rm low}=\frac{L\max(p_{z,+},p_{z,-})}{p_T},\qquad
  q_{\rm high}=\frac{U\min(p_{z,+},p_{z,-})}{p_T}.
  \label{v571:azimuth-bounds}
\end{equation}
Both daughters pass if $q_{\rm low}\leq |\sin\phi_*|\leq q_{\rm high}$.
Writing $[q]_0^1=\min(1,\max(0,q))$ and $[a]_+=\max(0,a)$, the accepted
azimuthal fraction and the remaining integral are
\begin{equation}
  F_\phi(c_*)=\frac2\pi
    \left[\arcsin\bigl([q_{\rm high}]_0^1\bigr)-
          \arcsin\bigl([q_{\rm low}]_0^1\bigr)\right]_+,\qquad
  \mathcal A_{\rm band}=\int_{-1}^{1}\!dc_*\,w(c_*)F_\phi(c_*).
  \label{v571:azimuth-fraction}
\end{equation}
We set $F_\phi=0$ when either daughter is backward or the nonzero-gap aperture
cannot be reached.  This fraction describes geometry under the stated
assumptions.  A selected mass spectrum additionally depends on production,
beam exposure, energy and direction distributions, trigger, reconstruction
and analysis requirements; it is not a measurement of $\mathcal A(m,x)$.
